%%
%% This is file `sample-sigconf.tex',
%% generated with the docstrip utility.
%%
%% The original source files were:
%%
%% samples.dtx  (with options: `sigconf')
%%
%% IMPORTANT NOTICE:
%%
%% For the copyright see the source file.
%%
%% Any modified versions of this file must be renamed
%% with new filenames distinct from sample-sigconf.tex.
%%
%% For distribution of the original source see the terms
%% for copying and modification in the file samples.dtx.
%%
%% This generated file may be distributed as long as the
%% original source files, as listed above, are part of the
%% same distribution. (The sources need not necessarily be
%% in the same archive or directory.)
%%
%%
%% Commands for TeXCount
%TC:macro \cite [option:text,text]
%TC:macro \citep [option:text,text]
%TC:macro \citet [option:text,text]
%TC:envir table 0 1
%TC:envir table* 0 1
%TC:envir tabular [ignore] word
%TC:envir displaymath 0 word
%TC:envir math 0 word
%TC:envir comment 0 0
%%
%%
%% The first command in your LaTeX source must be the \documentclass
%% command.
%%
%% For submission and review of your manuscript please change the
%% command to \documentclass[manuscript, screen, review]{acmart}.
%%
%% When submitting camera ready or to TAPS, please change the command
%% to \documentclass[sigconf]{acmart} or whichever template is required
%% for your publication.
%%

\documentclass[acmsmall,screen, nonacm, anonymous]{acmart}

\usepackage{stmaryrd}
\newcommand{\parr}{\mathbin{\bindnasrepma}}
% \usepackage{amssymb}
\usepackage{mathpartir}

\newcommand*{\ONECOLUMN}{}

% !TEX root = ../main.tex

\usepackage{docmute}
\usepackage{hyperref}
\usepackage{proof}

\usepackage{mathtools}
\usepackage{tikz}
\usepackage{tikz-network}
\usepackage{xspace}
\usepackage{xargs}
\usepackage{fontenc}

%\usepackage{minted}

\usepackage{subcaption}
\usepackage{caption}

\usepackage{tikz-cd}
\usepackage{tikzit}
% to reduce tikz picture white space
\usepackage{adjustbox}

% for quotienting
\usepackage{xfrac} 

\usepackage{import}

% category names
%\newcommand{\MdaCospans}{{\catname{MACsp_{D}(Hyp_{\Sigma})}}}

\newcommand\HypI[1]{\textbf{HypI}(#1)}
\newcommand\MdaCospans{\textbf{MHypI}(\Sigma)}
\newcommand{\Ecospans}{{\catname{EHypI(\Sigma)}}}
\newcommand{\MdaEcospans}{{\catname{MEHypI(\Sigma)}}}
\newcommand{\WellTypedMdaEcospans}{{\catname{MEHypI(\Sigma)}}}
% conflict
\newcommand{\hashtag}{{\#}}
\newcommand{\consistency}{{\smile}}
% There is already a remark in the preamble, but it does
% not work for some reason
\newtheorem{remark}{Remark}


\usepackage{xifthen}
\ifthenelse{\isundefined{\ismain}}{
  \newcommand{\ismain}{0}
}{}

\usepackage{enumitem}
 
\newcommand{\ignora}[1]{ }

\ignora{
\usepackage{aliascnt}
% alias-counters
\newaliascnt{definition}{thm}
\newaliascnt{proposition}{thm} 
\newaliascnt{lemma}{thm}
\newaliascnt{corollary}{thm}
\newaliascnt{conjecture}{thm}
\newaliascnt{remark}{thm}
\newaliascnt{example}{thm}
\newaliascnt{examples}{thm}
\newaliascnt{assumption}{thm}
\newaliascnt{construction}{thm}
\newaliascnt{claim}{thm}

% \theoremstyle{plain}
\theoremstyle{definition}

\newtheorem{theorem}[thm]{Theorem}
\newtheorem{proposition}[proposition]{Proposition} 
\aliascntresetthe{proposition}
\newtheorem{lemma}[lemma]{Lemma}
\aliascntresetthe{lemma}
\newtheorem{corollary}[corollary]{Corollary}
\aliascntresetthe{corollary}
\newtheorem{conjecture}[conjecture]{Conjecture}
\aliascntresetthe{conjecture}
\newtheorem{remark}[remark]{Remark}
\aliascntresetthe{remark}
\newtheorem{assumption}[assumption]{Assumption}
\aliascntresetthe{assumption}
\newtheorem{construction}[construction]{Construction}
\aliascntresetthe{construction}
\newtheorem{claim}[claim]{Claim}
\aliascntresetthe{claim}

\theoremstyle{definition}
\newtheorem{definition}[definition]{Definition}
\aliascntresetthe{definition}
\newtheorem{example}[example]{Example}
\aliascntresetthe{example}
\newtheorem{examples}[examples]{Examples}
\aliascntresetthe{examples}
}

\usepackage{cleveref} % after hyperref!
% possibly new thm environment
% \newtheorem{satz}[thm]{Satz}


% \newcommand{\TODO}[1]{{{\color{red}[\textbf{TODO:} #1]}}}
% \newcommand{\TODO}[1]{}

% !TEX root = ../main.tex

\newcommand{\missing}[1]{{\color{red}\bfseries [TODO]}}

\newcommand{\egg}{\texorpdfstring{\MakeLowercase{\texttt{egg}}}{\texttt{egg}}\xspace}
\newcommand{\Egg}{\texorpdfstring{\MakeLowercase{\texttt{egg}}}{\texttt{egg}}\xspace}
% \newcommand{\egg}{Lego\xspace}
% \newcommand{\Egg}{Lego\xspace}
\newcommand{\egraphs}{\mbox{e-graphs}\xspace}
\newcommand{\egraph}{\mbox{e-graph}\xspace}
\newcommand{\Egraph}{\mbox{E-graph}\xspace}
\newcommand{\Egraphs}{\mbox{E-graphs}\xspace}
\newcommand{\eclass}{\mbox{e-class}\xspace}
\newcommand{\Eclass}{\mbox{E-class}\xspace}
\newcommand{\enode}{\mbox{e-node}\xspace}
\newcommand{\eclasses}{\mbox{e-classes}\xspace}
\newcommand{\enodes}{\mbox{e-nodes}\xspace}
\newcommand{\Enodes}{\mbox{E-nodes}\xspace}
% \newcommand{\regraph}{Regraph\xspace}
% \newcommand{\Regraph}{Regraph\xspace}
\newcommand{\sz}{Szalinski\xspace}
\newcommand{\find}{\texttt{find}\xspace}

\newcommand{\equivid}{\equiv_{\sf id}}
\newcommand{\equivnode}{\equiv_{\sf node}}
\newcommand{\equivterm}{\equiv_{\sf term}}

\newcommand{\congrinv}{$\mathcal{I}_c$\xspace}

\newcommand{\egglogo}[1][]{\protect\includegraphics[height=1em, #1]{egg.png} }
\newcommand{\eggurl}{\url{https://github.com/mwillsey/egg}}

% TODOs
% https://tex.stackexchange.com/questions/9796/how-to-add-todo-notes

\newcommand{\eggtodo}[1]{\textcolor{Red}{\textbf{[CITE: #1]}}}

\newcommand{\defTodo}[2]{%
  \expandafter\newcommand\csname #1\endcsname[1]{%
    \todo[linecolor=#2,backgroundcolor=#2!25,bordercolor=#2,inline,size=\tiny]{\textbf{#1}: ##1}}}

\newcommand{\defTODO}[2]{%
  \expandafter\newcommand\csname #1\endcsname[1]{%
    \todo[linecolor=#2,backgroundcolor=#2!25,bordercolor=#2,inline,size=\tiny,caption={\textbf{(#1 LONG TODO)}}]{##1}}}

\newcommand{\AlekseiColor}{RoyalBlue}

% examples
% \newcommand{\RemyColor}{Red}
% \newcommand{\OliverColor}{OliveGreen}
% \newcommand{\ChandraColor}{Magenta}
% \newcommand{\PavelColor}{Cerulean}
% \newcommand{\ZachColor}{Plum}
% \newcommand{\BenColor}{Orchid}
% \newcommand{\JamesColor}{Salmon}


\defTodo{Aleksei}{\AlekseiColor}
% \defTodo{Max}{\MaxColor}

\defTODO{ALEKSEI}{\AlekseiColor}
% \defTODO{MAX}{\MaxColor}


% categorical stuff

\newcommand{\catname}[1]{\mathbf{#1}}

\newcommand{\Graph}{\catname{Graph}}
\newcommand{\Set}{\catname{Set}}

\newcommand\id{\textsf{id}}
\newcommand\sym{\textsf{sym}}

% tiny inline string diagrams
\newcommand{\comonoid}{%
	\begin{tikzpicture}[baseline=0.4ex, scale=0.4, line width=0.8pt]
		\path [use as bounding box] (0,0) rectangle (0.6,0.8);
		\draw (0.2,0) -- (0.2,0.4);
		\filldraw[black] (0.2,0.4) circle (0.08);
		\draw (0,0.8) .. controls (0,0.65) and (0.1,0.45) .. (0.2,0.4);
		\draw (0.4,0.8) .. controls (0.4,0.65) and (0.3,0.45) .. (0.2,0.4);
	\end{tikzpicture}%
}
\newcommand{\counit}{%
	\begin{tikzpicture}[baseline=0.4ex, scale=0.4, line width=0.8pt]
		\path [use as bounding box] (0,0) rectangle (0.4,0.8);
		\draw (0.2,0) -- (0.2,0.4);
		\filldraw[black] (0.2,0.4) circle (0.08);
	\end{tikzpicture}%
}

% camera-ready
\newcommand{\ifcameraready}[2]{\ifdefined\cameraready #2 \else #1 \fi}

\input{../sample.tikzstyles}
\input{../hypergraph.tikzstyles}
\input{../hypergraph.tikzdefs}


\definecolor{applegreen}{rgb}{0.55, 0.71, 0.0}
\definecolor{americanrose}{rgb}{1.0, 0.01, 0.24}
\definecolor{atomictangerine}{rgb}{1.0, 0.6, 0.4}
\definecolor{azure}{rgb}{0.0, 0.5, 1.0}

% These are for comments
\newcommand\question[1]{{\color{azure}#1}}
\newcommand\update[1]{{\color{americanrose}#1}}


\usepackage{minted}
\usepackage{csquotes}

% type derivations
\usepackage{mathpartir}


\makeatletter
\newsavebox{\@brx}
\newcommand{\llangle}[1][]{\savebox{\@brx}{\(\m@th{#1\langle}\)}%
  \mathopen{\copy\@brx\kern-0.5\wd\@brx\usebox{\@brx}}}
\newcommand{\rrangle}[1][]{\savebox{\@brx}{\(\m@th{#1\rangle}\)}%
  \mathclose{\copy\@brx\kern-0.5\wd\@brx\usebox{\@brx}}}
\makeatother


%%
%% \BibTeX command to typeset BibTeX logo in the docs
\AtBeginDocument{%
  \providecommand\BibTeX{{%
    Bib\TeX}}}


\settopmatter{printacmref=false}
\authorsaddresses{}

\begin{document}

\title{Notes on categorical logic and term calculi}

\author{Anonymous}


\maketitle

\section{Linear logic}


\textit{Linear logic} frequently pops up in the context of term calculus and categorical semantics and in fact is an internal logic of closed symmetric monoidal categories (with varied additional structures).
This section will briefly introduce what linear logic is and sketch its categorical semantics.
The corresponding term calculus is system L or $\overline{\lambda}\mu\widetilde{\mu}$-calculus.

Semantics is given by a symmetric monoidal category with enough tensors to support logical connectives.

\begin{center}
\begin{tabular}{c|cc}
{}& conjunction & disjunction\\
\hline
multiplicative & $\otimes$ & $\parr$\\
additive & $\&$ & $\oplus$\\
\end{tabular}
\end{center}

The difference between additive and multiplicative connectives is in how they treat contexts, \textit{e.g.}, multiplicative connectives split context into disjoint bits.
Additive conjunction is modelled by a Cartesian product and the multiplicative one is given by a tensor.
The modality $\!$ is interpreted by a comonad. 
Then we need to require that the comonad interacts with the tensor.


\section{Differential lambda calculus by ~\citet{DiffLambdaCalc}}

This is the calculus such that the model of the lambda terms are differentiable functions between vectors spaces.
The calculus is closely related to the notion of linearity and linear substitution.
In ordinary $\lambda$-calculus, application $(s) t$ substitutes $t$ into $s$.
In differential $\lambda$-calculus, we have two kinds of application:
\begin{itemize}

\item Nonlinear application (ordinary application) --- can duplicate arguments.

\item Linear application --- written as $D_{k} t \cdot u$ (term $t$ differentiated with respect to its k-th argument linearly applied to $u$)

\end{itemize}
The derivative $D t \cdot u$ says: \enquote{The linear effect on the result of $s$ when we change its input in the direction $u$.}
This is analogous to a directional derivative in calculus --- except `change' is in a syntactic/linear resource sense.

\begin{example}
Consider a term $(x) x$.
Its derivative (applied to $u$) $D (x) x \cdot u$ is $(u) x + x (u)$ (if we ignore linearity issues) where $+$ is a nondeterministic choice.
That is, the change is either propagated to the first $x$ or to the second.
\end{example}

The paper then proceeds to formally define $D$, $\partial x$ and the interaction of these with the $\beta$-reduction.
The calculus along with the reduction rules is shown to be strongly normalizing and confluent --- that is, equality of terms is straightforwardly decidable.

\section{Term graphs}

\textit{Term graphs} are essentially abstract syntax trees with sharing.
The consequence of this is that unlike our hypergraphs, they have to be rooted.
They share the same vocabulary with hypergraphs with some minor tweaks.
\begin{definition}
  A path between vertices $v_0$ and $v_n$ is a sequence $\langle v_0, i_1, v_1, \ldots, i_n, v_n \rangle$ of nodes and positive integers such that for $j = 1 \ldots n$, if $e$ is the unique edge such that $res(e) = v_{j-1}$, then $v_j$ is $i^{\text{th}}_j$ input to $e$.
\end{definition}

\begin{definition}
  A (hyper-) graph $G$ is a \textit{term graph} if
  \begin{enumerate}
    \item there is a vertex $root_{G}$
    \item $G$ is acyclic
    \item each node is the result of a unique edge (i.e. every edge represents a functional symbol with co-artity of 1)
  \end{enumerate}
\end{definition}

There is a canonical representative of an isomorphism class of term graphs given by a particular numbering of the vertices.
Given a vertex $v$ an \textit{access} path to $v$ is a sequence $\langle i_1, \ldots, i_n \rangle$ such that there is a path $\langle v_0, i_1, v_1, \ldots, i_n, v_n \rangle$ where $v_0 = root_G$ and $v_n = v$.
The set of access paths of $v$ is $Acc(v)$.
The numbering is then given by
\begin{enumerate}
  \item $v = Acc(v)$
  \item $e = res(e)$
\end{enumerate}

\textcolor{cyan}{We can do the same canonicalization for our (e-) hypergraphs by topologically sorting them. 
This gives a simple procedure for deciding equality modulo SMC equations.}

\[
\tikzfig{./figures/canonical_term_graph_example}
\]

\begin{remark}
  Several canonicalized term graphs can represent the same term (these differ in sharing/de-sharing of subtrees).
  Graphs ordered by sharing/de-sharing relation form a lattice where the least element (with no sharing) is denoted $\Delta G$ and the maximal element (with everything shared) is denoted by $\nabla G$.
\end{remark}

People then study the relation between term rewriting and term-graph rewriting.
For example, they answer the question of if term rewriting is confluent is the corresponding term-graph rewriting confluent?
What is interesting is that soundness and completeness are formulated the other way around.
Soundness is a statement of the form
\begin{proposition}
  For all term graphs $G$ and $H$, $G \xRightarrow{} H$ implies $term(G) \xrightarrow{+} term(H)$.
\end{proposition}

The other direction is then completeness.
To get completeness we need to rewrite term-graphs modulo bisimilarity (collapsing and copying).

\section{Categorical Type Theory}

This is the study of the relation between \textit{algebraic type theory} and categories with finite products.

\begin{definition}

An algebraic signature $\Sigma$ is given by:
\begin{itemize}

\item a collection of types;
\item a collection of functional symbols with given arity (possibly 0) but the co-arity is 1
\item a typing for each functional symbol $f$ which is a list of types $[\alpha_1, \ldots, \alpha_n, \alpha]$ such that $f : \alpha_1, \ldots, \alpha_{n} \to \alpha$
\end{itemize}

The raw terms are then given by
\[
 M \Coloneqq x \;|\; k \;|\; f(M_1, \ldots, M_{n})
\]
where $x$ is a variable from a countably infinite set of variables and $k$ is a functional symbol of arity 0.
\end{definition}

Given an algebraic theory $\Sigma$, an \textit{equation-in-context} is $\Gamma \vdash M = M' : \alpha$.

\begin{definition}

An algebraic theory is a pair $(\Sigma, \mathcal{A})$ where $\Sigma$ is a signature and $\mathcal{A}$ is a set of equations in context called the set of axioms.
\end{definition}

The theorems of a theory can then be derived from the axioms using reflexivity, symmetry, transitivity, substitution and context operations.
Categorical semantics then assigns objects to types and morphisms to type judgements.
In particular, substitution becomes morphism composition as (in the case of one variable)
\[
\inferrule*[right=Subst]
  { x : \alpha \vdash M : \beta  \\ y : \beta \vdash N : \gamma }
  { x : \alpha \vdash N[M / y] : \gamma }
\]

\[
\llbracket x : \alpha \vdash N[M / y] : \beta \rrbracket =  \llbracket y : \beta \vdash N : \gamma \rrbracket \circ \llbracket x : \alpha \vdash M : \beta \rrbracket
\]

Formally, letting 
\begin{align*}
\llbracket \alpha \rrbracket \in \text{obj}(C)\\
\llbracket k : \alpha \rrbracket = 1 \to \llbracket \alpha \rrbracket\\
\llbracket f : \alpha_1, \ldots, \alpha_n \to \alpha  \rrbracket = \llbracket \alpha_1 \rrbracket \otimes \ldots \llbracket \alpha_n \rrbracket \to \llbracket \alpha \rrbracket
\end{align*}

we get the lemma 

\begin{lemma}
\[
\llbracket \Gamma \vdash N[\vec{M} / \vec{x}] : \beta \rrbracket = \llbracket \Gamma \vdash N : \beta \rrbracket \circ \langle \llbracket \Gamma \vdash M_{1} : \alpha_{1} \rrbracket, \ldots, \llbracket \Gamma \vdash M_{n} : \alpha_{n} \rrbracket \rangle
\]

\end{lemma}

The assignment of $\llbracket - \rrbracket$ above is called a structure.
A structure is then \textit{models} the theory if all equations-in-context are equal morphisms.
The interpretation is sound if the structure also satisfies all the theorems.

\begin{definition}
  Given two models of a theory $Th$ in a category $\catname{C}$ we can define a homomorphism between two models such that diagrams for the interpretations commute.
  This gives rise to a category $Mod(Th,\catname{C})$.  
\end{definition}

\begin{definition}
  $\mathcal{FP}(C,D)$ for categories with products $C,D$ has as objects product-preserving functors and as morphisms it has natural transformations.
\end{definition}

The following claim is then can be made that for a given algebraic theory $Th$ we can find a category $Cl(Th)$ such that
\[
\mathcal{FP}(Cl(Th),D) \cong Mod(Th,D)
\]
That is, a model can be regarded as a functor. 
The equivalence above means there are functors between the two categories such that the composition of the functors are naturally isomorphic to identity functor.

Given a product preserving functor $F : C \to D$ we can lift it to $F_{*} : Mod(Th, C) \to Mod(Th,D)$.
Let $M$ be an object of $Mod(Th,C)$.
$F_{*}(M)$ is defined by putting $\llbracket \alpha \rrbracket_{F_{*}, M} = F\llbracket \alpha \rrbracket_{M}$ where $\alpha$ is the type of $Th$ and letting
$\llbracket f \rrbracket_{F_{*},M}$ to be given by composition
\[
  \llbracket \alpha_{1} \rrbracket_{F_{*}, M} \times \ldots \times \llbracket \alpha_{n} \rrbracket_{F_{*}, M} \cong F(\llbracket \alpha_{1} \rrbracket \times \ldots \times \llbracket \alpha_{n} \rrbracket) \xrightarrow{F(\llbracket f \rrbracket_{M})} F(\llbracket \beta \rrbracket_{M})
\]
For a morphism $h : M \to N$ in $Mod(Th,C)$, $F_{*}(h) = Fh$.

If $\phi : F \to G$ is a natural transformation between product preserving functors, we can define a homomorphism
\[
\phi_{*}M : FM \to GM\\
\phi_{\llbracket \alpha \rrbracket}M : F\llbracket \alpha \rrbracket \to G\llbracket \alpha \rrbracket_{M}
\]
of models.
Then we will define a family of modelling functors
\[
\mathcal{A}_{p_{M}} : \mathcal{FP}(C,D) \to Mod(Th,D)
\]
by setting $\mathcal{A}_{p_{M}}(F) = F_{*}(M)$ and $\mathcal{A}_{p_{M}}(\phi) = \phi_{*}M$.
A category (with finite products) $Cl(Th)$ is called a \emph{classifying} category of $Th$ if there is a model $G$ of $Th$ in $Cl(Th)$ such that for any other category $D$ the modelling functor
\[
\mathcal{A}_{p_{G}} : \mathcal{FP}(Cl(Th),D) \to Mod(Th,D)
\]
is an equivalence.
Any such $Cl(Th)$ is unique up to isomorphism.


\section{Categorical logic (as per Michael Shulman)}

The idea is to construct a free structure (with some nice syntactic properties, not every free structure will do) and then prove properties about non-free structures using the free structure and the homomorphism from this free structure.
For example, one could show that a property that holds for all groups by showing that it holds for a free group.

\begin{example}

Let a conjugation of $h$ by $g$, $h^{g}$ by $ghg^{-1}$.
For any $h,g,k$ we have $(hk)^{g} = h^{g}k^{g}$.

This statement can be proved by constructing a free group over a set of 3 elements $F_{3}$ and then by tracing homomorphisms to any other group.
\end{example}


\begin{example}

Another example is a type theory of posets (of pre-orders actually) that is a syntactical representation of morphisms in $\mathbf{Poset}$.
We can formulate the following question, is it true that for any objects $A,B,C,D,E$ satisfying

\[
A \leq B \qquad A \leq C \qquad D \leq A \qquad B \leq E \qquad D \leq C~,
\]
it is the case that $D \leq E$.

This can be answered by consulting the adjunction $\mathcal{F} \vdash U : \mathbf{Poset} \to \mathbf{RelGr}$ where the latter is a category of sets with binary relations.
Given a set $X$ in $\mathbf{RelGr}$, $\mathcal{F}X$ is a free poset on this set (i.e. its reflexive and transitive closure).
We can then say that $\mathcal{F}X$ has the same objects as $X$ and the morphisms are the judgements derivable from $X$ in a type theory of posets.
That is, given any other poset $A$ and a morphism $P : X \to UA$ there is a morphism $\mathcal{F}X \to A$ whose action on objects is given by $P$ and we can show by induction on the structure of the derivation that such morphism preserves the order.
\end{example}

\begin{example}

Another example stems from the adjunction $\mathcal{F} \vdash U : \mathbf{Cat} \to \mathbf{Graph}$ where the latter is the category of directed graphs.
To relate morphisms and derivations we will need to label the derivations.
For example, given two edges $f$ and $g$ in $\catname{Graph}(A,B)$, we get two primitive derivations

\[
\inferrule*[right=$f$]{\hspace{1em}}{A \vdash B} \qquad \inferrule*[right=$g$]{\hspace{1em}}{A \vdash B}
\]

Then we will need rules for composition, \textit{e.g.},

\[
\inferrule*[right=$\circ$]{
  \inferrule*{}{
    \inferrule*[right=$h$]{\;}{B \vdash C}
  }{
    \inferrule*[right=$g$]{\;}{A \vdash B}
  }
}{A \vdash C}
\]

Then, to make these derivations into proper categorical morphisms we need to quotient them by $\equiv$ (that is by associativity of $\circ$ and unitality of identity morphisms).
The claim is then that these data define a free category $\mathcal{F}G$ on a directed graph $G$.
The approach above (with quotienting) is being referred to as \textbf{cut-ful type theory}.

Another approach is to build the composition into axioms, \textit{i.e.}, for each $f \in \catname{Graph}(A,B)$ we have
\[
\inferrule*[right=$f$]{X \vdash A}{X \vdash B}
\]
Intuitively, the above rule is a post-composition with $f$.
This approach is called \textbf{cut-free type theory}.
We then can have a theorem that
\begin{theorem}
  Given a derivation $A \vdash B$ and $B \vdash C$ we can construct a derivation $A \vdash C$.
\end{theorem}
The theorem essentially states that cut-rule is admissible in this type theory.
The inference rules can be labelled with terms as follows (for $f : A \to B$)
\[
\inferrule*[right=$id$]{\;}{id_{A} : A \vdash A} \qquad \inferrule*[right=$f$]{\phi : X \vdash A}{f \circ \phi : A \vdash B}
\]
and we can identify terms with type derivations
\[
\inferrule*[right=$h$]{
\inferrule*[right=$g$]{
  \inferrule*[right=$f$]{
    \inferrule*[right=$id$]{\;}{id_{A} : A \vdash A}
  }{f \circ (id_{A}) : A \vdash B}
}{g \circ (f \circ id_{A}) : A \vdash C}
}{h \circ (g \circ (f \circ id_{A})) : A \vdash D}
\]

The syntax can be improved by introducing \textit{formal} variables for the left-hand sides of $\vdash$, so that the above term can be written as
\[
x : A \vdash h(g(f(x))) : D
\]
The identity and generator rules can now be written as
\[
\inferrule*[right=$id$]{\;}{x : A \vdash x : A} \qquad \inferrule*[right=$f$]{x : X \vdash M : A \hspace{2em} f \in \mathcal{G}(A,B)}{x : X \vdash f(M) : B}
\]
We can prove the cut admissibility using this notation as follows.
Given $x : A \vdash M : B$ and $y : B \vdash N : C$ we need to construct $x : A \vdash K : C$.
\begin{itemize}
  \item if $N = y$, then $C = B$ and $x : A \vdash M : B$ is the required term.
  \item if $N = f(N')$ for $N' : D$ and $f : D \to C$ we get a term $A \to D$ using the inductive hypothesis and after applying $f$ to it we get the required $A \to C$.
\end{itemize}

We can then say that a free category $\mathcal{F}(G)$ for a directed graph $G$ has the same objects and the morphisms are given by derivations in the type theory with composition given by substitution.
Substitution being associative and having unitality implies this is indeed a category and we can then define a functor $\mathcal{F}(G) \to \mathcal{A}$ for any category $\mathcal{A}$ by induction on the derivation using the morphism $G \to \mathcal{A}$.
This way no quotienting is required.
\end{example}

When we move to more complex theories, \textit{e.g.} the ones that have more complicated objects, such as meets in semilattices, we need to utilize categories with products.
Then, for the derivations we have two options (just like we had for a composition) regarding the projections $A \times B \to A$ and $A \times B \to B$.
\[
\inferrule*{A \vdash C}{A \times B \vdash C} \qquad \inferrule*{C \vdash A \times B}{C \vdash A}
\]

where the first is essentially precomposing a projection $\pi_1$ to $f: A \to C$ : $A \times B \xrightarrow{\pi_1} A \xrightarrow{f} C$.
The other one is postcomposition for $g : C \to A \times B$: $C \xrightarrow{g} A \times B \xrightarrow{\pi_1} A$.
Utilizing the first rule gives rise to \textit{sequent calculus} style type theories while the second one gives rise to \textit{natural deduction}.

In the case of a meet semilattice these rules get instantiated like this

\[
\inferrule*[right=$\land L_{1}$]{A \vdash C \quad \vdash B\;\text{type}}{A \land B \vdash C} \qquad \inferrule*[right=$\land L_{2}$]{B \vdash C \quad \vdash A\;\text{type}}{A \land B \vdash C} \qquad \inferrule*[right=$\land R$]{A \vdash B \quad A \vdash C}{A \vdash B \land C}
\]

and 

\[
\inferrule*[right=$E_{1}$]{X \vdash A \land B}{X \vdash A} \qquad \inferrule*[right=$E_{2}$]{X \vdash A \land B}{X \vdash B} \qquad \inferrule*[right=$\land I$]{X \vdash B \quad X \vdash C}{X \vdash B \land C}
\]

Such construction applies to (unary) type theories for categories with products and co-products where the version with variables is more naturally presented in natural deduction style.

\begin{example}

An unary type theory for a category with coproducts has the following inference rules.
For objects

\[
\inferrule*[]{X \in \mathcal{G}}{\vdash X\;\text{type}} \quad \inferrule*[]{\;}{\vdash \mathbf{0}\; \text{type}} \quad \inferrule*[]{\vdash A\;\text{type} \quad \vdash B\;\text{type}}{\vdash A + B\;\text{type}}
\]

For morphisms
\[
\inferrule*[right=$id$]{\vdash X\;\text{type}}{x : X \vdash x : X} \quad \inferrule*[right=$fI$]{f \in \mathcal{G}(A,B) \quad x : X \vdash M : A}{x : X \vdash f(M) : B}
\]
\[
\inferrule*[right=$\mathbf{0}E$]{x : X \vdash M : \mathbf{0} \quad \vdash C\;\text{type}}{x : X \vdash \text{match}_{0}(M) : C}
\]
\[
\inferrule*[right=$+I1$]{x : X \vdash M : A}{x : X \vdash inl(M) : A + B} \quad \inferrule*[right=$+I2$]{x : X \vdash M : B}{x : X \vdash inr(M) : A + B}
\]
\[
\inferrule*[right=$+E$]{x : X \vdash M : A + B \quad u : A \vdash P : C \quad v : B \vdash Q : C}{x : X \vdash \text{match}_{A+B}(M,u.P,v.Q) : C}
\]
noting that in the very last rule $v$ and $u$ are bound in $u.P$ and $v.Q$.
\end{example}

To enforce the universal property of the coproduct we need to introduce the following equations
\[
\text{match}(inl(y), u.P, v.Q) \equiv P[inl(y) / u]
\]
\[
\text{match}(inr(y), u.P, v.Q) \equiv Q[inl(y) / v]
\]
and
\[
\text{match}(y, u.P[inl(u)/y], v.P[inr(v)/y]) \equiv P
\]
where $P : A + B \to C$ and the rule basically says that pattern matching is redundant here.
And finally the rule (uniqueness of the arrow from the initial object)
\[
\text{match}_{\mathbf{0}}(y) \equiv P, \forall P
\]
\textcolor{red}{I am not sure what the process of coming up with such rules should be. Do I start with introduction/elimination rules and then try and fix the issues that break universal properties?}

\subsection{Simple Type Theories}
To reason about operations with multiple inputs we need to have multiple entities in the antecedent.
The right categorical setting to do this is \textit{multicategories}.

\begin{definition}

A multigraph $\mathcal{G}$ consists of a set of objects $\mathcal{G}_{0}$ together with for each object $B$ and every list of objects $(A_1,\ldots,A_n)$ a set of arrows $\mathcal{G}(A_1,\ldots,A_n;B)$.
\end{definition}

\begin{definition}
A multicategory $\mathcal{M}$ is a multigraph with the following structure and properties
\begin{itemize}
  \item For each object $A$, an identify arrow $id_{A} \in \mathcal{M}(A,A)$
  \item For any object $C$ and a list of objects $(B_1, \ldots, B_m)$ and $(A_{i1}, \ldots, A_{i n_{i}})$ for $1 \leq i \leq m$, a composition operation
  \begin{align*}
    \mathcal{M}(B_1, \ldots, B_m; C) \times \prod_{i=1}^{m}\mathcal{M}(A_{i1}, \ldots, A_{in_{i}};B_{i}) &\to \mathcal{M}(A_{11}, \ldots, A_{m n_{m}};C)\\
    (g, (f_1, \ldots, f_m)) &\mapsto g \circ (f_1, \ldots, f_m)
  \end{align*}
  \item For any $f \in \mathcal{M}(A_1, \ldots, A_n; B)$ we have
    \[
      id_{B} \circ (f) = f \qquad f \circ (id_{A_1}, \ldots, id_{A_n}) = f
    \] 
  \item For any $h, g_i, f_{ij}$ we have
    \[
      (h \circ (g_1, \ldots, g_m)) \circ (f_{11}, \ldots, f_{mn_m}) = h \circ (g_1 \circ (f_{11}, \ldots, f_{1 n_{1}}), \ldots, g_{m} \circ (f_{m 1}, \ldots, f_{m n_{m}}))
    \]
\end{itemize}
\end{definition}

A good example is the category of vector spaces and multilinear maps.
We can also introduce an indexed composition $g \circ_i f_i$ as a surag for $g \circ (f_1, \ldots, f_n)$ when all $f_j$ for $i \not = j$ are identities and $g \in \mathcal{M}(A_1, \ldots, A_n; B)$.
A bit more formally we have
\[
\circ_{i} : \mathcal{M}(B_1, \ldots, B_n; C) \times \mathcal{M}(A_1, \ldots, A_m) \to \mathcal{M}(B_1, \ldots, B_{i-1}, A_{1}, \ldots, A_{m}, B_{i + 1}, \ldots, B_n; C)
\]
\[
\tikzfig{./figures/single_comp}
\]
that satisfy the following properties
\begin{itemize}
  \item $id_{B} \circ_{1} f = f$
  \item $f \circ_{i} id_{B_{i}} = f$
  \item If $h$ is $n$-ary, $g$ is $m$-ary, and $f$ is $k$-ary, then
  \[
    (h \circ_i g) \circ_j f = 
  \begin{cases}
    (h \circ_j f) \circ_{i + k - 1} g \qquad &\text{if}\; j < i\\
    (h \circ_i (g \circ_{j - i + 1} f)) \qquad &\text{if}\; i \leq j < i + m\\
    (h \circ_{j-m+1} f) \circ_i g \qquad &j \geq i + m
  \end{cases}
  \]
\end{itemize}

\begin{definition}

A tensor product in $\mathcal{M}$ is an object $\bigotimes_{i}A_{i}$ together with a morphism $\chi \in \mathcal{M}(A_1,\ldots,A_n; \bigotimes_{i}A_{i})$ such that all $(- \circ_{i} \chi_{i})$ are bijections.
\[
\mathcal{M}(B_{1},\ldots,B_{n},A_1,\ldots,A_m,C_{1},\ldots,C_{k};D) \xrightarrow{\sim} \mathcal{M}(B_{1},\ldots,B_{n},\bigotimes_{i}A_i,C_{1},\ldots,C_{k};D)
\]
A unit $\mathbf{1}$ is given by a nullary tensor product.
\end{definition}

We call a multicategory representable if it is equipped with a chosed unit object and for every pair of objects with their tensor.

Here are some extensions with additional structures

\paragraph*{Symmetric monoidal structure} If $\mathcal{C}$ is a symmetric monoidal category, we have symmetry isomorphisms $A_1 \otimes \ldots \otimes A_n \xrightarrow{\sim} A_{\sigma 1} \otimes \ldots \otimes A_{\sigma n}$ for any permutation $\sigma \in S_{n}$.
Thus, by precomposing with these isomorphisms, we obtain functions between multicategorical hom-sets
\[
\sigma^* : \mathcal{C}(A_{\sigma 1}, \ldots, A_{\sigma n}; B) \to \mathcal{C}(A_1,\ldots,A_n;B)
\]
That is, suppose we have a multimorphism $f : A_{\sigma 1}, \ldots, A_{\sigma n} \to B$. We can precompose it with the isomorphism above to get
\[
\sigma^{*}(f) : A_1 \otimes \ldots \otimes A_n \xrightarrow{\sim} A_{\sigma 1} \otimes \ldots \otimes A_{\sigma n} \xrightarrow{f} B
\]
\textcolor{red}{I do not really see what is the point of these multicategories}
\textcolor{cyan}{My take on this is that is more general as it does not impose any structure on the list in the antecedent and one could make it into a tensor if they want but it also works if the multicategory is not representable.
It also does not require to invoke any associativity as the antecedent is a literal list.
Additional structure is imposed by universal properties.
}
\paragraph*{Cartesian structure} If $\mathcal{C}$ is a cartesian monoidal category, we have symmetries in addition to dual maps $A \to A \times A$ and projections $A \times B \to B$.
In general for any function $\sigma : \{1, \ldots, m\} \to \{1, \ldots, n\}$ we have a unique morphism
\[
A_1 \times \ldots \times A_n \to A_{\sigma 1} \times \ldots A_{\sigma m}
\]
whose component $A_1 \times \ldots \times A_n \to A_{\sigma k}$ is $\sigma k$-th projection.
Precomposition with those gives
\[
\sigma^{*} : \mathcal{C}(A_{\sigma 1}, \ldots, A_{\sigma m}; B) \to \mathcal{C}(A_{1}, \ldots, A_n; B)
\] 

\begin{definition}
  Let $\mathfrak{N}$ be a full subcategory of $\catname{Set}$ whose objects are sets ${1, \ldots, n}$ for all integers $n \geq 0$.
  We regard it as a cocartesian strict monoidal category, under the disjoint union operation $\{1, \ldots, n\} \sqcup \{1,\ldots,m\} = \{1, \ldots, n + m\}$.
  Moreover, for any $\sigma : \{1, \ldots, m\} \to \{1, \ldots, n\}$ and $k_1, \ldots, k_n$, let $\sigma \wr (k_1, \ldots, k_n)$ denote the composite function
  \[
  \{1, \ldots, \sum_{i=1}^{m}k_{\sigma i}\} \xrightarrow{\sim} \bigsqcup_{i=1}^{m}\{1, \ldots, k_{\sigma i}\} \xrightarrow{\hat{\sigma}} \bigsqcup_{j=1}^{n}\{1, \ldots, k_j\} \xrightarrow{\sim} \{1, \ldots, \sum_{j=1}^{n} k_j\}
  \]
  where $\hat{\sigma}$ acts as an identity from the $i$-th summand to the $(\sigma i)^{\text{th}}$ summand.
  A \textbf{faithful cartesian club} is a category $\mathfrak{S} \subseteq \mathfrak{N}$ such that 
  \begin{enumerate}
    \item $\mathfrak{S}$ contains all the objects of $\mathfrak{N}$.
    \item $\mathfrak{S}$ is closed under the cocartesian structure, i.e. if $\sigma$ and $\tau$ are morphisms of $\mathfrak{S}$ then so is $\sigma \sqcup \tau$.
    \item $\mathfrak{S}$ is closed under $\wr$, i.e. whenever it contains $\sigma$ it also contains $\sigma \wr (k_1, \ldots, k_n)$
  \end{enumerate}
\end{definition}

The strange-looking $\sigma \wr (k_1, \ldots, k_n)$. 
Here's what's happening: Start with a function $\sigma : \{1, \ldots, m\} \to \{1, \ldots, n\}$ which says that each $i$ gets assigned one of the $n$ slots.
Now, instead of just $n$ slots, suppose each slot $j$ is \enquote{blown up} into $k_j$ copies.
The construction $\sigma \wr (k_1, \ldots, k_n)$ gives a new function
\[
\{1, \ldots, \sum_{i=1}^{m}k_{\sigma i}\} \to \{1, \ldots, \sum_{j=1}^{n}k_j\}
\]
which is what you get when you first group the inputs according to $\sigma$, and then send them into the right block of the \enquote{expanded} target.


We get symmetric monoidal category when $\mathfrak{S}$ consists of all bijections and cartesian structure when it consists of all functions.

\begin{definition}
  Let $\mathfrak{S}$ be a faithful cartesian club. An $\mathfrak{S}$-multicategory is a multicategory $\mathcal{M}$ together with operations
  \[
  \mathcal{M}(A_{\sigma 1}, \ldots, A_{\sigma m};B) \to \mathcal{M}(A_1, \ldots, A_n;B)
  \]
  \[
  f \mapsto f \circ \sigma^{*}
  \]
  for all functions $\sigma : \{1, \ldots, m\} \to \{1, \ldots, n\}$ in $\mathfrak{S}$, satisfying the following axioms:
  \begin{enumerate}
    \item $f\sigma^{*}\tau^{*} = f(\tau \sigma)^{*}$
    \item $f(id_n) = f$
    \item $g \circ (f_1\sigma^{*}_{1}, \ldots, f_{n}\sigma^{*}_n) = (g \circ (f_1, \ldots, f_n))(\sigma_1 \sqcup \ldots \sqcup)^{*}$
    \item $g \sigma^{*} \circ (f_1, \ldots, f_n) = (g \circ (f_{\sigma 1}, \ldots, f_{\sigma_m}))(\sigma \wr (k_1, \ldots, k_n))^{*}$, where $k_i$ is arity of $f_i$.
  \end{enumerate}
\end{definition}

Finally, we observe that closedness can be naturally characterized multicategorically.
Suppose for simplicity that $\mathfrak{S}$ contains at least all bijections. 
Then we say an $\mathfrak{S}$-multicategory is closed if for each pair of objects $A$ and $B$ there is a
specified object $A \multimap B$ and a morphism $\chi : A \multimap B \to B $ postcomposition
with which defines bijections
\[
(\chi \circ -) \mathcal{M}(C_1, \ldots, C_n; A \multimap B) \to \mathcal{M}(C_1, \ldots, C_n, A; B)
\]

\subsection{Type theory for (closed) symmetric monoidal categories}

With all of this in hand we can define type theory for symmetric monoidal closed categories (in natural deduction style).
Below $\text{Shuf}(\Gamma_1, \ldots, \Gamma_n; \Phi)$ says that $\Phi$ is a shuffle of $\Gamma_1, \ldots, \Gamma_n$, i.e. $\Phi$ is obtained by rifle-shuffling.


\[
\inferrule*[right=$id$]{\vdash A\;\text{type}}{x : A \vdash x: A}
\]
\[
\inferrule*[right=$fI$]{f \in \mathcal{G}(A_1, \ldots, A_n; B)\;\Gamma_1 \vdash M_1 : A_1, \ldots, \Gamma_n \vdash M_n : A_n\;\text{Shuf}(\Gamma_1,\ldots,\Gamma_n;\Phi)}{\Phi \vdash f(M_1, \ldots, M_n) : B}
\]
\[
\inferrule*[right=$\otimes I$]{\Gamma \vdash N : A \qquad \Delta \vdash M : B \qquad \text{Shuf}(\Gamma,\Delta;\Phi)}{\Phi \vdash \{N \otimes M\} : A \otimes B}
\]
\textcolor{red}{The rule above seems to imply the tensor is not strict.}
\textcolor{red}{Can the below rule be made admissible? It looks like a composition.}
\[
\inferrule*[right=$\otimes E$]{\Psi \vdash M : A \otimes B \quad \Gamma, x : A, y : B \vdash N : C \quad \text{Shuf}(\Gamma,\Psi;\Phi)}{\Phi \vdash \text{match}_{A \otimes B}(M, xy.N) : C}
\]
\[
\inferrule*[right=$\mathbf{1}I$]{\;}{() \vdash \star : \mathbf{1}} \qquad \inferrule*[right=$\mathbf{1}E$]{\Psi \vdash M : \mathbf{1} \qquad \Gamma,\Delta \vdash N : A}{\Gamma, \Psi. \Delta \vdash \text{match}_{\mathbf{1}}(M,N) : A}
\]
\[
\inferrule*[right=$\multimap I$]{\Gamma, x : A \vdash M : B}{\Gamma \vdash \lambda x . M : A \multimap B}
\]
\textcolor{red}{Should there be a symmetry involved in the rule above?}
\[
\inferrule*[right=$\multimap E$]{\Gamma \vdash M : A \multimap B \quad \Delta \vdash N : A \quad \text{Shuf}(\Gamma, \Delta; \Phi)}{\Phi \vdash M N : B}
\]

To implement the universal properties, we need the following equations for the tensor ($\beta$ and $\eta$ laws)
\[
\text{match}_{\otimes}(\{M,N\}, xy.P) \equiv P[M / x, N /y]
\]
\[
\text{match}_{\otimes}(M, xy.N[\{x,y\}/ u]) \equiv N[M / u]
\]

\[
\text{match}_{\mathbf{1}}(\star, N) \equiv N \qquad \text{match}_{\mathbf{1}}(M,N[\star / u]) \equiv N[M / u]
\]

The above rules enjoy the properties of exchange being admissible, substitution being admissible, terms representing derivations.
In particular, substitution is associative and has an interchange property.
The initiality theorem is then that a free closed symmetric monoidal category on a multigraph is given by this type theory.
\textcolor{cyan}{What we might want to do is augment the above rules for the type theory to generate a semilattice enriched symmetric monoidal category and formulate the corresponding initiality theorem.
Terms calculus may provide an interesting angle on the differences between cartesian and non-cartesian case, for example in the former we may disregard some inputs inside one of the components of $f_1 + \ldots + f_n$ while in non-cartesian case we can not and therefore all the variables used by $f_i$ must be used by $f_j$ for $i \not = j$.}

\begin{remark}
  A term uniquely determines a derivation only when paired with its context.
  For example, below are two different derivations (in the presence of weakening).
  \[
  \inferrule*[]{
    \inferrule*[]{
      \vdash A\;\text{type}\;(x : A) \in \Gamma
      }
      {x : A, z : C, y : B \vdash x : A
      }
     \quad 
      \inferrule*[]{
        \vdash B\;\text{type}\;(y : B) \in \Gamma
        }
        {x : A, z : C, y : B \vdash y : B
        }
        }
      {x : A, z : C, y : B \vdash \{x,y\} : A \otimes B}
  \]

  \[
  \inferrule*[]{
    \inferrule*[]{
      \vdash A\;\text{type}\;(x : A) \in \Gamma
      }
      {x : A, y : B \vdash x : A
      }
     \quad 
      \inferrule*[]{
        \vdash B\;\text{type}\;(y : B) \in \Gamma
        }
        {x : A, y : B \vdash y : B
        }
        }
        {x : A, y : B \vdash \{x,y\} : A \otimes B}
  \]
\end{remark}

So far we can say that the context correspond to the input interface of our hypergraphs and the pushout along the interfaces correspond to a substitution.
It is unclear how to make a claim that a subgraph correspond to a term as they are defined using different structure.

The bottom line is that the above syntax is a \enquote{nice} syntax for morphisms in symmetric monoidal categories.
For example, consider the below morphism, first represented as a string diagram, then represented as a raw term and finally represented using the syntax above.

\[
\tikzfig{./figures/string-diagram-example}
\]

\[
(f \otimes g);\sigma_{D,E};h
\]
which is not unique as it is equal e.g. to the following term
\[
(f;id_{D} \otimes g);\sigma_{D,E};\sigma_{E,D};\sigma_{D,E};h
\]

\[
\inferrule*{
  \inferrule*{
    \inferrule*{f \in \mathcal{G} \quad x : A \vdash x : A \quad y : B \vdash y : B}{x : A, y : B \vdash f(x,y) : D}
    \qquad
    \inferrule*{g \in \mathcal{G} \quad z : C \vdash z : C}{z : C \vdash g(z) : E}
  }{x : A, y : B, z : C \vdash \{f(x,y), g(z)\} : D \otimes E}
  \qquad
  \inferrule*{h \in \mathcal{G} \quad v : E \vdash v : E \quad u : D \vdash u : D}
  {u : D, v : E \vdash h(v,u)}
}{x : A, y : B, z : C \vdash \text{match}(\{f(x,y), g(z)\},uv.h(v,u))}
\]


\begin{figure}
\begin{align*}
x[M/x] &= M\\
y[M/x] &= y\\
f(N_1, \ldots, N_n)[M/x] &= f(N_1, \ldots, N_i[M/x], \ldots, N_n) \quad \text{if $x$ occurs in $N_i$}\\
\{N_1 \otimes N_2\}[M/x] &= \{N_1[M/x] \otimes N_2\} \quad \text{if $x$ occurs in $N_1$}\\
\{N_1 \otimes N_2\}[M/x] &= \{N_1 \otimes N_2[M/x]\} \quad \text{if $x$ occurs in $N_2$}\\
\text{match}(N_1,xy.N_2)[M/x] &= \text{match}(N_1[M/x], xy.N_2) \quad \text{if $x$ occurs in $N_1$}\\
\text{match}(N_1,xy.N_2)[M/x] &= \text{match}(N_1, xy.N_2[M/x]) \quad \text{if $x$ occurs in $N_2$}\\
(\lambda y . N)[M / x] &= \lambda y . N[M/x]\\
(N_1 N_2)[M / x] &= (N_1[M/x])N_2 \quad \text{if $x$ occurs in $N_1$}\\
(N_1 N_2)[M / x] &= (N_1)N_2[M/x] \quad \text{if $x$ occurs in $N_2$}
\end{align*}
\caption{Substitution definition}
\label{fig:sub-def}
\end{figure}

\begin{lemma}[Substitution (Cut rule) is admissible]
  Given derivations $\Gamma \vdash M : A$ and $\Delta, x : A, \Phi \vdash N : B$ we can derive
  \[
  \Gamma, \Delta, \Phi \vdash N[M / x]
  \]
  Substitution is defined in \cref{fig:sub-def}.
  Note, in particular, that since contexts are disjoint (every variable is used once) we only recurse in a single subterm.
\end{lemma}
\begin{proof}
  We will induct on the derivation of $\Delta, x : A, \Phi \vdash N : B$.
  \begin{itemize}
    \item If it is the $id$ rule then $\Delta = \varnothing = \Phi$ and $N : B = x : A$ and $B = A$. Then we have $M[x/x] = M$.
    \item  
    \[
    \inferrule*{f \in \mathcal{G} \quad \Gamma_1 \vdash N_1, \ldots, \Gamma_n \vdash N_n}{\Delta , x : A, \Phi \vdash f(N_1, \ldots, N_n) : B = N : B}
    \]
    Then it must be that there exists $1 \leq i \leq n$ such that $x \in \Gamma_i$ and we can apply inductive hypothesis to $\Gamma \vdash M : A$ and $\Gamma', x : A, \Gamma'' \vdash N_i : A_i$ to obtain $N_i[M/x] : A_i$.
    By applying the same derivation rule again we get $f(N_1, \ldots, N_i[M/x], \ldots, N_n)$.
    \item The other rules are analogous.
  \end{itemize}
\end{proof}

With the lemma above we can prove, for example, that $A \otimes \mathbf{1}$ is isomorphic to $A$.
For this we need to construct two morphisms $A \otimes \mathbf{1} \to A$ and $A \to A \otimes \mathbf{1}$ and show that their composition is identity.
First, let's build a morphism from $A \to A \otimes \mathbf{1}$

\[
\inferrule*{
  x : A \vdash x : A \qquad () \vdash \star : \mathbf{1}
}{x : A \vdash \{x,\star\} : A \otimes \mathbf{1}}
\]
the morphism $A \otimes \mathbf{1} \to A$ is then
\[
\inferrule*{
\inferrule*{}{z : A \otimes \mathbf{1} \vdash z : A \otimes \mathbf{1}}
\qquad
\inferrule*{y : \mathbf{1} \vdash y : \mathbf{1} \qquad x : A \vdash x : A}{x : A, y : \mathbf{1} \vdash \text{match}_{\mathbf{1}}(y,x) : A}
}{z : A \otimes \mathbf{1} \vdash \text{match}_{A \otimes \mathbf{1}}(z, xy.\text{match}_{\mathbf{1}}(y,x)) : A}
\]
Then we have
\begin{itemize}
  \item \[
  \inferrule*{
  \inferrule*{
  \inferrule*{        
  \inferrule*{
            \inferrule*{
  w : A \vdash w : A \qquad () \vdash \star : \mathbf{1}
}{w : A \vdash \{w,\star\} : A \otimes \mathbf{1}}
\qquad
\inferrule*{
\inferrule*{}{z : A \otimes \mathbf{1} \vdash z : A \otimes \mathbf{1}}
\qquad
\inferrule*{y : \mathbf{1} \vdash y : \mathbf{1} \qquad x : A \vdash x : A}{x : A, y : \mathbf{1} \vdash \text{match}_{\mathbf{1}}(y,x) : A}
}{z : A \otimes \mathbf{1} \vdash \text{match}_{A \otimes \mathbf{1}}(z, xy.\text{match}_{\mathbf{1}}(y,x)) : A}
}{w : A \vdash \text{match}_{A \otimes \mathbf{1}}(z, xy.\text{match}_{\mathbf{1}}(y,x))[\{w, \star\}/z]}}{
w : A \vdash \text{match}_{A \otimes \mathbf{1}}(\{w, \star\}, xy.\text{match}_{\mathbf{1}}(y,x))}}{
w : A \vdash \text{match}_{\mathbf{1}}(y,x)[w / x,  \star / y]}}{
w : A \vdash \text{match}_{\mathbf{1}}(\star,w) \equiv w
}
      \]
  \item The other direction is substitution 
  \[
  \inferrule*{
  \inferrule*{z : A \otimes \mathbf{1} \vdash \{w, \star\}[\text{match}_{A \otimes \mathbf{1}}(z, xy.\text{match}_{\mathbf{1}}(y,x)) / w]}
{z : A \otimes \mathbf{1} \vdash \{\text{match}_{A \otimes \mathbf{1}}(z, xy.\text{match}_{\mathbf{1}}(y,x)), \star\}}}
{z : A \otimes \mathbf{1} \vdash \text{\textcolor{red}{What's next?}}}
  \]
  
\end{itemize}

\begin{lemma}[Substitution is associative and interchanging]
  \begin{enumerate}
  \item If $\Gamma \vdash M : A$ and $\Delta, x : A, \Delta' \vdash N : B$ and $\Psi, y : B, \Psi' \vdash P : C$, then
  \[
  (P[M/x])[N/y] = P[M[N/y]/x]
  \]
  \item If $\Gamma \vdash M : A$ and $\Delta \vdash N : B$ and $\Psi, x : A, \Psi', y : B, \Psi'' \vdash P : C$, then
  \[
  (P[M/x][N/y]) = P[N/y][M/x]
  \]
  \end{enumerate}
\end{lemma}

\begin{proposition}
  A free (closed) symmetric monoidal category on a (hyper-) multigraph $\mathcal{G}$ is given by the type theory above.
  The objects are types and the morphisms are derivations (terms).
\end{proposition}

\subsection{Semilattice-enriched symmetric monoidal categories}

We want to extend the rules of the previous section with

\[
\inferrule*{\Gamma \vdash M_1 : A, \ldots, \Gamma \vdash M_n : A}{
  \Gamma \vdash M_1 + \ldots + M_n
}
\]

$M_1, \ldots, M_n$ look different depending on what structural rules we have.
In the most general case they all should use the same variables, but in a cartesian case they can ignore some of them.

and prove the analogous statements, i.e. admissibility of substitution, associativity of substitution etc.
Then we can say that a free semilattice-enriched symmetric monoidal category is given by appropriate type theory.

\begin{example}
  Recall the raw morphism representation for our e-graph example
  \[(a \otimes ((2;\comonoid) \otimes (1;id));\sym);((* \otimes \counit + (id \otimes \counit \otimes id);<\!\!<) \otimes id);/~,\]
  We can render it using the  \enquote{enhanced} term syntax.
  \[
  \inferrule*{
  \inferrule*{
  \inferrule*{\vdash a : Int \quad \vdash 2 : Int}{\vdash a * 2}
  \quad
  \inferrule*{\vdash a : Int \quad \vdash 1 : Int}{\vdash a <\!\!< 1}
  }{\vdash (a * 2) + (a <\!\!< 1)}
  \quad
  \inferrule*{}{\vdash 1 : Int}
  }
  {((a * 2) + (a <\!\!< 1)) / 2}
  \]
\end{example}

\bibliographystyle{ACM-Reference-Format}
\bibliography{../bibliography}

\end{document}
